\chapter{Domain-Specific WARP Graphs (Physics, Biology,
Logic)}\label{domain-specific-rmgs-physics-biology-logic}

\begin{quote}
When Every Discipline Is Just a Rewrite System in Drag
\end{quote}

Up to this point we've treated WARP graphs as a universal substrate --- the
underlying fabric from which computation, geometry, and physics all
emerge. But the real magic happens when you zoom into specific domains.
Suddenly the abstract becomes concrete, and every scientific field
reveals itself as a special case of the same ontological engine.

This chapter gives the reader a hard jolt of recognition: physics is a
domain-specific WARP graph, biology is a domain-specific WARP graph, logic is a
domain-specific WARP graph, and so is everything else.

No exceptions.

Once you see the world this way, you can't unsee it.

\begin{center}\rule{0.5\linewidth}{0.5pt}\end{center}

\subsection{1. What Makes a Domain-Specific
WARP Graph?}\label{what-makes-a-domain-specific-rmg}

A Domain-Specific WARP Graph (DS-WARP) is:

\begin{itemize}
\tightlist
\item
  a WARP graph whose rewrite rules encode the ``laws'' of a given domain,
\item
  whose constraints enforce the domain's invariants,
\item
  whose curvature reflects the domain's geometry,
\item
  and whose stable subgraphs represent the phenomena that domain
  studies.
\end{itemize}

Physics = specific rule families. Biology = specific rule families with
selective pressures. Logic = specific rewrite constraints defining
provability.

And each DS-WARP is a patch on the grand MRMW manifold.

The same underlying ontology. Different local physics.

\begin{center}\rule{0.5\linewidth}{0.5pt}\end{center}

\subsection{2. DS-WARP \#1: Physics}\label{ds-rmg-1-physics}

The Original Rewrite Machine

The universe is, to first approximation, a WARP graph where:

\begin{itemize}
\tightlist
\item
  nodes are physical configurations
\item
  edges are allowable transitions
\item
  DPO rules encode local causal interactions
\item
  curvature arises from constraint propagation
\item
  stable subgraphs are conserved quantities
\item
  rewrite bundles produce quantum superposition
\item
  and measurement collapses rewrite branches into a single worldline
\end{itemize}

This is not metaphor. It is isomorphism.

Classical mechanics emerges when rewrite curvature is low and rules are
deterministic. Quantum mechanics emerges when rewrite bundles are wide
and constraints collapse them. General relativity emerges when curvature
is non-uniform across the manifold. Thermodynamics emerges from the
combinatorics of rewrite congestion.

Physics is the canonical DS-WARP, and the physical laws we observe are
simply the rule set of our local cosmic patch.

Rewriting is physics. Physics is rewriting.

The rest of the sciences are echoes.

\begin{center}\rule{0.5\linewidth}{0.5pt}\end{center}

\subsection{3. DS-WARP \#2: Biology}\label{ds-rmg-2-biology}

Life as Rewrite Systems Under Selective Pressure

Biology is where WARP graphs get weird\emd{}in a good way.

If physics is the base WARP graph, biology is a higher-order WARP graph riding on top
of it. Its rules include:

\begin{itemize}
\tightlist
\item
  replication
\item
  mutation
\item
  selection
\item
  self-reference
\item
  homeostasis
\item
  emergent constraints
\item
  information-bearing structures
\end{itemize}

These are rewrites acting on rewrites. A recursion on the WARP toolkit
itself.

A gene is a stable subgraph. A cell is a rewrite-bounded pocket
universe. An organism is a massively parallel rewrite manifold
maintaining constraints. A population is a multiversal bundle of nearby
worldlines under divergent selective curvature.

Evolution is literally:

\begin{lstlisting}
\text{Selection Pressure} = -\nabla_R \Phi_{\text{fitness}}
\end{lstlisting}

The environment defines the potential field. Organisms follow gradient
descent. Successful subgraphs survive.

Biology is physics with feedback. DS-WARP \#2 exposes that life is not an
exception but an intensification of universal rewrite law.

\begin{center}\rule{0.5\linewidth}{0.5pt}\end{center}

\subsection{4. DS-WARP \#3: Logic}\label{ds-rmg-3-logic}

Proof Systems as Rewrite Constraint Machines

Logic is the cleanest DS-WARP in terms of structure:

\begin{itemize}
\tightlist
\item
  propositions are nodes
\item
  inference rules are rewrites
\item
  proofs are worldlines
\item
  consistency is curvature neutrality
\item
  contradictions are negative curvature collapse points
\item
  semantics are stability conditions
\item
  non-termination = infinite worldlines
\item
  normal forms are stable attractors
\end{itemize}

A logical proof is a path in a rewrite universe whose objective is to
reach a stable subgraph representing a theorem.

\paragraph{Gödel's incompleteness?}\label{guxf6dels-incompleteness}

A region of the WARP graph where local curvature cannot be flattened globally.

\paragraph{Curry-Howard?}\label{curry-howard}

A direct isomorphism between two DS-WARP graphs (logic and computation) with
different constraints.

\paragraph{Modal logics?}\label{modal-logics}

WARP graphs with extra curvature channels.

\paragraph{Probabilistic logics?}\label{probabilistic-logics}

WARP graphs with width-weighted rewrite bundles.

Logic, seen properly, is not about truth\emd{}it's about allowed rewrites.

\begin{center}\rule{0.5\linewidth}{0.5pt}\end{center}

\subsection{5. Four Other DS-WARP Graphs Worth
Naming}\label{four-other-ds-rmgs-worth-naming}

To reveal the generality of the frame

We will expand these later in the formal appendix:

\begin{itemize}
\tightlist
\item
  Economics → agents as rewrite regions, markets as constraint fields
\item
  Neuroscience → synaptic rewrites under plasticity curvature
\item
  Sociology → multi-agent WARP graphs with high-dimensional constraint layers
\item
  Machine Learning → parameter manifolds as rewrite-influenced curvature
  surfaces
\end{itemize}

Every discipline has a rewrite ontology. Most just haven't admitted it
yet.

\begin{center}\rule{0.5\linewidth}{0.5pt}\end{center}

\subsection{6. Domain Interfaces: The Crucial
Insight}\label{domain-interfaces-the-crucial-insight}

The real power emerges when DS-WARP graphs are composable via spans.

A physics WARP graph and a biology WARP graph can interface through energy/matter
constraints. A biology WARP graph and a logic WARP graph can interface through
cognitive structure. A machine learning WARP graph and a logic WARP graph can
interface through invariants. A physics WARP graph and a computation WARP graph can
interface through measurement.

This is the missing meta-scientific glue. Without a unified substrate,
you can't fuse different fields formally.

WARP graphs solve that.

There is only one mathematics under all these domains: graphs, rules,
constraints, curvature, and worldlines.

Everything else is dialect.

\begin{center}\rule{0.5\linewidth}{0.5pt}\end{center}

\subsection{7. Why DS-WARP Graphs Matter for
CΩMPUTER}\label{why-ds-rmgs-matter-for-cux3c9mputer}

The runtime doesn't operate in some abstract nowhere. It will be called
to simulate domains:

\begin{itemize}
\tightlist
\item
  physics simulations
\item
  biological processes
\item
  logical inference engines
\item
  cognitive models
\item
  agent-based systems
\item
  formal program behavior
\end{itemize}

Without DS-WARP graphs, the runtime is universal but useless. With DS-WARP graphs, the
runtime becomes a general-purpose engine for simulating reality, natural
and artificial.

This is where CΩMPUTER transitions from theory to applied
omnidisciplinarity.
